%% ============================================================
%% QCCK 第三章 中文版
%% 格式参考：Overleaf 项目 6aa379e8fec9886e0960f9e9（article + ctex 单栏中文稿）
%% 内容来源：ch3/第三章修改版1.md（用户中文底稿，逐句原样保留）
%% 处理：去掉「第N句：」写作脚手架；公式排成 equation 环境；
%%       「图 X」引用删除（本版不收录图）；[REF] 落为 bracewell 引用
%% 编译：xelatex（ctex 需要 xelatex 或 lualatex）
%% ============================================================

% [已改造] 本文件现为 main.tex 的分章正文（导言区已并入 main.tex），由 main 统一编译。


\section{QCCK：量子跨变量耦合核}

本章提出量子跨变量耦合核 QCCK，将变量间耦合显式地构造出来。
首先，频域对齐得到去除时移与尺度差异的对齐频谱；随后，QCCE 将变量信息编码为纠缠量子态；
最后，量子核比较各变量量子态的重叠程度，得到可解释的耦合矩阵 $K$，供下游主干在此基础上进行跨变量建模与预测。

\subsection{QCCK 总览}

QCCK 的整体结构由频域对齐、QCCE 与量子核三个环节构成，对每个样本显式度量变量两两之间的耦合，输出可解释的耦合矩阵 $K\in\mathbb{R}^{V\times V}$。
频域对齐作为 QCCK 的输入处理环节，由谱变换与时间轴、频率轴两级对齐构成，将每个变量的原始序列变换为对齐频谱 $S$，消除变量间的时移与周期尺度等环境差异，使耦合度量只刻画变量间的结构关系。对齐频谱 $S$ 是 QCCE 中相位通道的基础。
QCCE 作为 QCCK 方法的核心，通过振幅通道与相位通道将每个变量的信息编码为量子态，并借助参数化量子线路在比特间建立纠缠，构造出能反映数据变量特征的纠缠量子态 $|\psi_v\rangle$，作为量子核度量变量间耦合的输入。
量子核作为 QCCK 的耦合度量环节，以 QCCE 得到的各变量量子态为输入，按式~(\ref{eq:qkernel}) 的保真度核计算任意两个变量 $i$、$j$ 的重叠度作为耦合强度，即 $K[i,j] = |\langle\psi_i|\psi_j\rangle|^2$，其中 $\langle\cdot|\cdot\rangle$ 为 $\mathbb{C}^{32}$ 上的标准内积，$|\psi_i\rangle$、$|\psi_j\rangle$ 分别为变量 $i$、$j$ 经 QCCE 编码得到的量子态。$K\in\mathbb{R}^{V\times V}$ 取值在 $[0,1]$，其每个元素表示一对变量的耦合强度。
因此，$K$ 的可解释性体现在它把变量间的耦合落成一份可检查、可比较的矩阵：据此能定位强耦合的变量对，也能看出单个变量与其余各变量的耦合分布；用热力图呈现会更直观，且每个样本对应一份 $K$，可观察耦合结构随样本的变化。
综上所述，针对现有双向状态空间类主干以扫描方式隐式传递变量间关系、导致变量间耦合难以被显式刻画与利用的问题，我们提出了 QCCK。它显式度量变量间的耦合，得到具备可解释性的耦合矩阵 $K$，把原本隐式的变量间依赖变为可读、可解释的依据。

\subsection{频域对齐}

针对不同变量的波形普遍存在时移与主导周期差异、若不消除会使结构相似的变量被误判为互不相关的问题，为使耦合度量只反映变量间的结构关系，我们提出频域对齐：对每个变量的原始序列进行谱分析，并在时间轴与频率轴上分别消除时移与周期尺度差异，得到对齐频谱 $S$。
对齐后，时移与主导周期的差异不再干扰后续的耦合度量，结构相似而时移、周期不同的变量不会再被误判为互不相关。
记 $V$ 个变量的输入序列为 $x_1,\dots,x_V$，其中 $x_v \in \mathbb{R}^L$ 表示变量 $v$ 在连续 $L$ 个时刻上的原始观测序列，也就是该变量随时间变化的一段数据。这 $V$ 条序列取自同一段连续时刻，共同构成一个样本。我们对每个这样的 $x_v$ 施加实数快速傅里叶变换 rFFT，得到复谱 $X_v(f) \in \mathbb{C}^F$，其中 $F = \lfloor L/2 \rfloor+1$ 为频点个数，$f = 0,1,\dots,F-1$ 为频率索引。
我们定义了幅度谱与相位谱，如公式~(\ref{eq:amp-spec})、(\ref{eq:phase-spec}) 所示，二者分别刻画各频率分量的强度与时间位置。

\textbf{幅度谱。}
\begin{equation}
\label{eq:amp-spec}
A_v(f) = |X_v(f)|.
\end{equation}
$X_v(f)$ 是复数，对它取模 $|X_v(f)|$ 就得到频率 $f$ 分量的幅值，$A_v(f)$ 越大说明该频率成分越强，因此幅度谱刻画各频率分量的强度。当主导周期不同、波形在时间上伸缩时，幅度谱的主峰位置会随之移动，所以尺度差异体现在幅度谱上，频率轴对齐将依据幅度谱的主峰进行。

\textbf{相位谱。}
\begin{equation}
\label{eq:phase-spec}
\phi_v(f) = \arg X_v(f).
\end{equation}
$\arg$ 取复数的辐角，得到频率 $f$ 分量的相位，它决定该分量在时间轴上的起始位置，因此相位谱刻画各频率分量的时间位置。当不同变量的波形存在整体时移时，波形平移改变各分量在时间轴上的起始位置，即改变它们的相位，因此时移差异体现在相位谱上，时间轴对齐将针对相位谱进行。
令 $\hat{f}_{\mathrm{peak}}(v) = \arg\max_{f\geq 1} A_v(f)$ 为幅度谱主频，对应主导周期。忽略直流分量 $f=0$，因为直流只反映序列均值，不对应某个周期。
由尺度定理~\cite{bracewell2000fourier}，时间尺度差异表现为幅度谱~(\ref{eq:amp-spec}) 在频率轴上的缩放，据此归一化频率坐标，如公式~(\ref{eq:freq-align}) 所示：
\begin{equation}
\label{eq:freq-align}
\tilde{f} = f / \hat{f}_{\mathrm{peak}},
\end{equation}
使所有变量的主峰对齐至同一频率位置，消除主导周期差异。
时移校正需要一个共同的时序基准，我们以该样本内除变量 $v$ 外其余变量的平均序列为基准，使基准不受变量 $v$ 自身取值主导。令 $\hat{\delta}_v$ 表示变量 $v$ 相对该基准的时移估计，取为变量 $v$ 与基准互相关峰值对应的滞后占序列长度的比例。
由时移定理~\cite{bracewell2000fourier}，时移差异表现为相位谱~(\ref{eq:phase-spec}) 中随频率线性增长的相位项，据此校正相位谱，如公式~(\ref{eq:time-align}) 所示：
\begin{equation}
\label{eq:time-align}
\phi'_v(f) = \phi_v(f) - 2\pi f\hat{\delta}_v,
\end{equation}
使各变量校正后的相位谱对齐至同一时间原点，仅保留波形自身的结构信息，消除变量间的时移差异。
我们在对齐后的统一频率坐标上取固定的采样区间 $\tilde{f}\in[0,2]$，该区间覆盖主峰 $\tilde{f}=1$ 至其两倍频率以内的谐波结构；在该区间上取 $M = 32$ 个均匀采样点并作线性插值，得重采样幅度谱 $\tilde{A}_v \in \mathbb{R}^M$ 与重采样相位谱 $\tilde{\phi}_v \in \mathbb{R}^M$，拼接为对齐频谱，如公式~(\ref{eq:spectrum}) 所示。
\begin{equation}
\label{eq:spectrum}
S_v = [\,\tilde{A}_v\,;\,\tilde{\phi}_v\,] \in \mathbb{R}^{2M} = \mathbb{R}^{64}.
\end{equation}
由时移定理与尺度定理可知，对齐频谱 $S_v$ 对时移与时间尺度变换保持不变，维度恒为 $2M$，与输入长度 $L$ 无关。
综上所述，频域对齐把每个变量的原始序列变换为对齐频谱 $S_v$，为 QCCE 的相位通道提供同基准的频域输入。

\subsection{QCCE 编码}

编码环节 QCCE 是 QCCK 的重要组成部分，负责把每个变量的信息编码成量子态。它以主干输出的隐状态 $h_v$ 为振幅通道的输入、以频域对齐得到的对齐频谱 $S_v$ 为相位通道的输入，经参数化量子线路在比特间建立纠缠，得到反映变量特征的纠缠量子态 $|\psi_v\rangle$，作为量子核度量变量间耦合的输入。
变量 $v$ 的编码输入有两路：主干隐状态 $h_v\in\mathbb{R}^d$ 与对齐频谱 $S_v\in\mathbb{R}^{2M}$。双向状态空间时序主干把每个变量的整段历史观测编码为 $d$ 维表示并逐层更新，取 QCCK 所在编码层之后变量 $v$ 的表示即得 $h_v$，$h_v$ 刻画该变量自身的时序特征。对齐频谱由频域对齐从变量 $v$ 的原始序列得到，携带各频率成分对齐后的强度与相位信息。
为解决两类异质输入要进入同一个量子态且互不干扰的问题，我们定义了振幅通道与相位通道。
振幅通道把变量自身特征编码为各计算基态的概率：它以 $h_v$ 为输入，设振幅投影矩阵 $W_r\in\mathbb{R}^{32\times d}$，令
\begin{equation}
\label{eq:amp-channel}
p(h_v) = \mathrm{softmax}(W_r\cdot h_v) \in \mathbb{R}^{32}.
\end{equation}
$p(h_v)$ 各分量非负且和为 1，给出各计算基态被选中的概率。
相位通道把对齐频谱编码为各基态的相位：它以 $S_v$ 为输入，设相位投影矩阵 $W_\theta\in\mathbb{R}^{32\times 2M}$，令
\begin{equation}
\label{eq:phase-channel}
\theta(S_v) = W_\theta\cdot S_v \in \mathbb{R}^{32}.
\end{equation}
$\theta(S_v)$ 各分量给出对应计算基态的相位。
两个通道由此各司其职：振幅通道给出的概率决定各基态的振幅，相位通道给出的相位决定各基态间的相对相位；由于振幅与相位是量子态相互独立的两个自由度，两类信息互不干扰，使合成的量子态能够同时携带变量自身特征与频域信息。振幅、相位两通道的投影矩阵以及参数化线路的旋转角均对每个变量独立可学习，公式为行文简洁省略变量下标。
记 $b\in\{0,1\}^5$ 为 5 个量子比特的基态标记，$|b\rangle$ 为对应计算基态，共 $2^5=32$ 个；$p_b(h_v)$、$\theta_b(S_v)$ 分别为 $p(h_v)$、$\theta(S_v)$ 的第 $b$ 个分量，给出基态 $|b\rangle$ 的概率与相位。我们把这 32 个基态按各自的概率与相位叠加，即得变量 $v$ 在进入量子线路前的编码态
\begin{equation}
\label{eq:encode}
|\chi_v\rangle = \sum_{b\in\{0,1\}^5} \sqrt{p_b(h_v)} \cdot e^{\,i\cdot\theta_b(S_v)} \cdot |b\rangle.
\end{equation}
至此，变量 $v$ 的编码态 $|\chi_v\rangle$ 构造完成，作为变量 $v$ 进入后续量子线路演化的输入。式~(\ref{eq:encode}) 是式~(\ref{eq:joint-encoding}) 的联合编码在本文中的具体形式：振幅由主干隐状态给出，相位由对齐频谱给出。
为使编码态在比特间引入纠缠、得到量子核度量耦合所用的纠缠量子态，我们让编码态 $|\chi_v\rangle$ 经参数化量子线路演化。每个变量沿各自独立可学习的线路演化，变量 $v$ 的线路酉记作 $U_v$，一层由 RY、RZ 旋转门与 CNOT 纠缠门构成，其算符形式为
\begin{equation}
\label{eq:circuit}
U_v = \mathrm{CNOT}_{4,5}\cdot\mathrm{CNOT}_{3,4}\cdot\mathrm{CNOT}_{2,3}\cdot\mathrm{CNOT}_{1,2} \cdot \Bigl(\bigotimes_{q=1}^{5} RZ_q(\beta_q)\cdot RY_q(\alpha_q)\Bigr),
\end{equation}
其中 CNOT 依次作用于相邻比特对 $(1,2)$、$(2,3)$、$(3,4)$、$(4,5)$；$\alpha_q$、$\beta_q$ 为变量 $v$ 线路的可学习旋转角。式中 $R_Y$、$R_Z$ 旋转门与 CNOT 纠缠门的定义见附录~\ref{app:quantum}。编码态连续经过两层相同的 $U_v$，得到变量 $v$ 的最终量子态
\begin{equation}
\label{eq:final-state}
|\psi_v\rangle = U_v\cdot U_v\cdot|\chi_v\rangle.
\end{equation}
$|\psi_v\rangle$ 是变量 $v$ 的量子态表示：主干隐状态 $h_v$ 经振幅通道写成各计算基态的概率，对齐频谱 $S_v$ 经相位通道写成基态间的相位，两层线路又在比特间引入纠缠。它把主干隐状态 $h_v$ 与对齐频谱 $S_v$ 这两类异质输入收进同一个单位向量，每个变量各得一个 $|\psi_v\rangle$、构造一致，是连接编码输入与后续耦合度量的枢纽。
本节完成 QCCK 的编码环节 QCCE：以主干隐状态与对齐频谱为输入，经振幅、相位两个通道分别承载变量自身特征与频域信息，再经参数化量子线路两层的演化，把每个变量编码为一个归一化的纠缠量子态 $|\psi_v\rangle$。下一步，量子核将基于这些量子态度量变量间的耦合。

\subsection{量子核}

量子核作为 QCCK 的耦合度量环节，负责度量变量两两之间的耦合强度。它以各变量经 QCCE 编码并演化得到的纠缠量子态 $|\psi_v\rangle$ 为输入，对样本中的任意两个变量计算两态的重叠程度，得到耦合矩阵 $K\in\mathbb{R}^{V\times V}$。
变量 $i$、$j$ 的耦合强度由二者的量子态共同决定。量子核按式~(\ref{eq:qkernel}) 计算两态的重叠度，即内积模长的平方 $|\langle\psi_i|\psi_j\rangle|^2$，其中 $\langle\cdot|\cdot\rangle$ 为 $\mathbb{C}^{32}$ 上的标准内积，$|\psi_i\rangle$、$|\psi_j\rangle$ 分别为变量 $i$、$j$ 经 QCCE 编码并演化得到的纠缠量子态。
两量子态的重叠度即二者的保真度，因此本文采用的耦合度量就是式~(\ref{eq:qkernel}) 给出的保真度核。
保真度核刻画两态的重叠程度：两态越接近，输出越接近 1；两态越趋于正交，输出越接近 0；
当两变量为同一变量时，量子态与自身完全重叠，输出恒为 1，所以耦合矩阵的对角元素恒为 1。
综上，QCCK 以原始序列与主干隐状态为输入，经频域对齐、QCCE 与量子核三个环节，显式度量变量间的耦合，得到耦合矩阵 $K\in\mathbb{R}^{V\times V}$。$K$ 中任意两个变量的耦合强度都能直接读取，配合热力图即可定位强耦合的变量对，耦合关系因此不会隐没于黑盒权重之中。
可解释性由此随度量一同产生，模型所依据的变量间关系都能通过该矩阵直接呈现与核对。
在此基础上，第四章进一步将 QCCK 以 QCCK Layer 的形式嵌入状态空间主干的编码器层之间，并利用耦合矩阵 $K$ 引导跨变量信息融合，使显式度量得到的耦合关系直接参与变量表示更新与预测。



